\documentclass[11pt,a4paper]{article}

% ───────────── PAQUETES ─────────────
\usepackage[utf8]{inputenc}
\usepackage[T1]{fontenc}
\usepackage[spanish]{babel}
\usepackage[margin=2.5cm]{geometry}
\usepackage{graphicx}
\usepackage{float}
\usepackage{longtable}
\usepackage{booktabs}
\usepackage{array}
\usepackage{xcolor}
\usepackage{hyperref}
\usepackage{listings}
\usepackage{enumitem}
\usepackage{fancyhdr}
\usepackage{titlesec}
\usepackage{amssymb}

% ───────────── COLORES ─────────────
\definecolor{codebg}{RGB}{245,245,245}
\definecolor{codegreen}{RGB}{0,128,0}
\definecolor{codegray}{RGB}{120,120,120}
\definecolor{codepurple}{RGB}{128,0,128}
\definecolor{primary}{RGB}{30,90,160}

% ───────────── HYPERREF ─────────────
\hypersetup{
    colorlinks=true,
    linkcolor=primary,
    urlcolor=primary,
    citecolor=primary,
    pdftitle={PetTask - README},
    pdfauthor={Equipo PetTask}
}

% ───────────── LISTINGS (código) ─────────────
\lstdefinestyle{phpstyle}{
    backgroundcolor=\color{codebg},
    commentstyle=\color{codegreen},
    keywordstyle=\color{codepurple}\bfseries,
    stringstyle=\color{orange},
    numberstyle=\tiny\color{codegray},
    basicstyle=\ttfamily\small,
    breaklines=true,
    captionpos=b,
    keepspaces=true,
    numbers=left,
    numbersep=8pt,
    showspaces=false,
    showstringspaces=false,
    showtabs=false,
    tabsize=2,
    frame=single,
    rulecolor=\color{codegray}
}
\lstset{style=phpstyle}

% ───────────── ENCABEZADO ─────────────
\pagestyle{fancy}
\fancyhf{}
\fancyhead[L]{\small\textcolor{primary}{\textbf{PetTask}}}
\fancyhead[R]{\small\textcolor{primary}{README}}
\fancyfoot[C]{\thepage}
\renewcommand{\headrulewidth}{0.4pt}

% ───────────── TÍTULOS ─────────────
\titleformat{\section}{\Large\bfseries\color{primary}}{\thesection}{1em}{}
\titleformat{\subsection}{\large\bfseries\color{primary}}{\thesubsection}{1em}{}

% ───────────── INICIO ─────────────
\begin{document}

% ───── PORTADA ─────
\begin{titlepage}
    \centering
    \vspace*{3cm}
    {\Huge\bfseries\color{primary} PetTask\par}
    \vspace{0.5cm}
    {\Large CRUD de Tareas con Mascotas\par}
    \vspace{0.3cm}
    {\large Laravel + Blade (sin base de datos)\par}
    \vspace{2cm}
    \rule{\linewidth}{0.5pt}\\[0.4cm]
    {\large\textbf{Documentación del Proyecto --- README}\par}
    \rule{\linewidth}{0.5pt}\\[2cm]

    \begin{tabular}{rl}
        \textbf{Asignatura:} & Desarrollo Web con Laravel \\
        \textbf{Profesora:}  & Yeniza \\
        \textbf{Equipo:}     & PetTask \\
    \end{tabular}

    \vfill
    {\large \today\par}
\end{titlepage}

\tableofcontents
\newpage

% ───── 1. DESCRIPCIÓN ─────
\section{Descripción del proyecto}

Aplicación web desarrollada en \textbf{Laravel} que permite gestionar tareas asignadas a mascotas usando los \textbf{cuatro métodos HTTP} (\texttt{GET}, \texttt{POST}, \texttt{PUT/PATCH} y \texttt{DELETE}). El proyecto \textbf{no utiliza base de datos}: toda la información se almacena temporalmente en la \textbf{sesión} de Laravel, cumpliendo con los requisitos de la asignación.

% ───── 2. EQUIPO ─────
\section{Integrantes del equipo}

\begin{table}[H]
\centering
\renewcommand{\arraystretch}{1.3}
\begin{tabular}{|l|l|}
\hline
\textbf{Integrante} & \textbf{Responsabilidad} \\ \hline
Emir      & Vista de detalle (\texttt{show.blade.php}) \\ \hline
Ángel     & Lógica del controlador (\texttt{PetTaskController.php}) \\ \hline
Alexander & Layout principal y listado (\texttt{app.blade.php} / \texttt{index.blade.php}) \\ \hline
Chofos    & Formularios (\texttt{create.blade.php} / \texttt{edit.blade.php}) \\ \hline
Angélica  & Documentación (\texttt{README.md} / \texttt{readme2.tex}) \\ \hline
Mariana   & Pruebas, capturas y organización de la demo \\ \hline
\end{tabular}
\caption{Distribución de tareas del equipo.}
\end{table}

% ───── 3. TECNOLOGÍAS ─────
\section{Tecnologías utilizadas}

\begin{itemize}[leftmargin=1.5em]
    \item \textbf{Laravel 10/11} --- framework PHP.
    \item \textbf{Blade} --- motor de plantillas.
    \item \textbf{Bootstrap 5} --- estilos de la interfaz.
    \item \textbf{Sesiones de Laravel} --- almacenamiento temporal (sin base de datos).
    \item \textbf{Git / GitHub} --- control de versiones.
\end{itemize}

% ───── 4. REQUISITOS ─────
\section{Requisitos previos}

\begin{itemize}[leftmargin=1.5em]
    \item \textbf{PHP} $\geq$ 8.1
    \item \textbf{Composer} instalado (\url{https://getcomposer.org/})
    \item \textbf{Node.js y npm} (opcional, solo si se compilan assets)
    \item \textbf{Git} instalado (\url{https://git-scm.com/})
\end{itemize}

% ───── 5. INSTALACIÓN ─────
\section{Instalación paso a paso}

Ejecutar los siguientes comandos en la terminal:

\begin{lstlisting}[language=bash, caption={Comandos de instalación}, numbers=none]
# 1. Clonar el repositorio
git clone https://github.com/HolaSoyNickName/PetTask.git

# 2. Entrar a la carpeta del proyecto
cd PetTask

# 3. Instalar dependencias de PHP
composer install

# 4. Copiar el archivo de variables de entorno
cp .env.example .env

# 5. Generar la clave de la aplicacion
php artisan key:generate

# 6. (Opcional) Instalar dependencias de frontend
npm install && npm run build

# 7. Levantar el servidor local
php artisan serve
\end{lstlisting}

Luego abrir el navegador en: \\
\centerline{\url{http://localhost:8000/tasks}}

\textit{\textbf{Nota:} este proyecto \textbf{no necesita configurar base de datos}. Los datos viven en la sesión y se reinician al cerrar el navegador o al volver a ejecutar \texttt{php artisan serve}.}

% ───── 6. ESTRUCTURA ─────
\section{Estructura del proyecto (resumen)}

\begin{lstlisting}[language=bash, numbers=none, caption={Estructura de carpetas}]
PetTask/
|-- app/Http/Controllers/
|   `-- PetTaskController.php     # Logica CRUD con sesion
|-- resources/views/
|   |-- layouts/app.blade.php     # Layout general (navbar)
|   `-- tasks/
|       |-- index.blade.php       # Listado de tareas
|       |-- create.blade.php      # Formulario de creacion
|       |-- edit.blade.php        # Formulario de edicion
|       `-- show.blade.php        # Detalle de tarea
|-- routes/
|   `-- web.php                   # Definicion de rutas
`-- README.md
\end{lstlisting}

% ───── 7. TABLA DE RUTAS ─────
\section{Tabla de rutas del sistema}

\begin{longtable}{|p{2cm}|p{2.8cm}|p{2.8cm}|p{2.5cm}|p{4.5cm}|}
\hline
\textbf{Método} & \textbf{URI} & \textbf{Acción} & \textbf{Ruta} & \textbf{Descripción} \\ \hline
\endhead
GET     & \texttt{/tasks}              & \texttt{index()}   & \texttt{tasks.index}   & Muestra la lista de todas las tareas \\ \hline
GET     & \texttt{/tasks/create}       & \texttt{create()}  & \texttt{tasks.create}  & Muestra el formulario para crear tarea \\ \hline
POST    & \texttt{/tasks}              & \texttt{store()}   & \texttt{tasks.store}   & Guarda una nueva tarea en la sesión \\ \hline
GET     & \texttt{/tasks/\{id\}}         & \texttt{show()}    & \texttt{tasks.show}    & Muestra el detalle de una tarea \\ \hline
GET     & \texttt{/tasks/\{id\}/edit}    & \texttt{edit()}    & \texttt{tasks.edit}    & Muestra el formulario de edición \\ \hline
PUT/PATCH & \texttt{/tasks/\{id\}}       & \texttt{update()}  & \texttt{tasks.update}  & Actualiza una tarea existente \\ \hline
DELETE  & \texttt{/tasks/\{id\}}         & \texttt{destroy()} & \texttt{tasks.destroy} & Elimina una tarea de la sesión \\ \hline
\caption{Tabla de rutas del CRUD.}
\end{longtable}

\noindent\textcolor{primary}{$\checkmark$} Los cuatro métodos HTTP quedan cubiertos: \textbf{GET} (index, create, show, edit), \textbf{POST} (store), \textbf{PUT/PATCH} (update) y \textbf{DELETE} (destroy).

% ───── 8. CAPTURAS ─────
\section{Capturas de pantalla}

Las capturas deben guardarse en la carpeta \texttt{docs/screenshots/} con los nombres indicados.

\subsection{Pantalla principal --- Listado de tareas (index)}
Aquí se muestran todas las mascotas con sus tareas y los botones \textbf{Ver}, \textbf{Editar} y \textbf{Eliminar}.

\begin{figure}[H]
    \centering
    % \includegraphics[width=0.85\linewidth]{docs/screenshots/01-index.png}
    \caption{Pantalla principal del sistema.}
\end{figure}

\subsection{Formulario de creación (create)}
Formulario para registrar una nueva tarea: título, mascota y descripción.

\begin{figure}[H]
    \centering
    % \includegraphics[width=0.85\linewidth]{docs/screenshots/02-create.png}
    \caption{Formulario de creación de tarea.}
\end{figure}

\subsection{Vista de detalle (show)}
Pantalla con la información completa de una tarea seleccionada.

\begin{figure}[H]
    \centering
    % \includegraphics[width=0.85\linewidth]{docs/screenshots/03-show.png}
    \caption{Detalle de una tarea.}
\end{figure}

\subsection{Formulario de edición (edit)}
Permite modificar los datos de una tarea existente.

\begin{figure}[H]
    \centering
    % \includegraphics[width=0.85\linewidth]{docs/screenshots/04-edit.png}
    \caption{Formulario de edición de tarea.}
\end{figure}

\subsection{Confirmación al eliminar (destroy)}
Mensaje de confirmación al borrar una tarea del listado.

\begin{figure}[H]
    \centering
    % \includegraphics[width=0.85\linewidth]{docs/screenshots/05-delete.png}
    \caption{Confirmación de eliminación.}
\end{figure}

% ───── 9. GUION DEMO ─────
\section{Guion para la demo en vivo (5 minutos)}

\begin{table}[H]
\centering
\renewcommand{\arraystretch}{1.25}
\begin{tabular}{|c|l|l|p{6cm}|}
\hline
\textbf{Tiempo} & \textbf{Quién} & \textbf{Qué muestra} \\ \hline
0:00 --- 0:30 & NickName  & Introducción: nombre del proyecto, equipo y objetivo \\ \hline
0:30 --- 1:30 & Alexander & Recorrido por la pantalla principal (\texttt{index}) y el menú \\ \hline
1:30 --- 2:30 & Chofos    & Demo de crear (\texttt{create}) y editar (\texttt{edit}) una tarea \\ \hline
2:30 --- 3:15 & Emir      & Ver el detalle (\texttt{show}) de una tarea \\ \hline
3:15 --- 4:00 & Ángel     & Explicación breve del controlador y la sesión \\ \hline
4:00 --- 4:45 & Angélica  & Recorrido del README y tabla de rutas \\ \hline
4:45 --- 5:00 & Todos     & Cierre y preguntas \\ \hline
\end{tabular}
\caption{Guion de la demostración en vivo.}
\end{table}

% ───── 10. LÓGICA DE SESIÓN ─────
\section{¿Cómo funciona el almacenamiento en sesión?}

En lugar de una base de datos, el controlador \texttt{PetTaskController} guarda un arreglo de tareas dentro de la sesión de Laravel:

\begin{lstlisting}[language=PHP, caption={Ejemplo simplificado del método \texttt{store()}}]
public function store(Request $request)
{
    $tasks = session('tasks', []);
    $tasks[] = [
        'id'          => uniqid(),
        'titulo'      => $request->titulo,
        'mascota'     => $request->mascota,
        'descripcion' => $request->descripcion,
    ];
    session(['tasks' => $tasks]);

    return redirect()->route('tasks.index');
}
\end{lstlisting}

Así se logra un \textbf{CRUD completo sin base de datos}, cumpliendo con los requisitos de la actividad.

% ───── 11. CHECKLIST ─────
\section{Checklist de entrega}

\begin{itemize}[leftmargin=1.5em]
    \item[$\boxtimes$] Repositorio público en GitHub.
    \item[$\boxtimes$] Commits de todos los integrantes.
    \item[$\boxtimes$] \texttt{README} con instrucciones, capturas y tabla de rutas.
    \item[$\boxtimes$] Rutas declaradas con los 4 métodos HTTP.
    \item[$\boxtimes$] Controlador con lógica de sesión.
    \item[$\boxtimes$] Vistas Blade (\texttt{index}, \texttt{create}, \texttt{show}, \texttt{edit}, layout).
    \item[$\boxtimes$] Demo de 5 minutos organizada.
\end{itemize}

% ───── 12. LICENCIA ─────
\section{Licencia}

Proyecto académico desarrollado para la asignatura impartida por la \textbf{Profesora Yeniza}. Uso educativo únicamente.

\vspace{1cm}
\begin{center}
\textit{Hecho con amor por el equipo \textbf{PetTask}}
\end{center}

\end{document}